
\section{\extrahard\ 非诺特环上的凝聚模}

本节是为这样两种人写的：可能要同非\trNoetherian{}环打交道的人，或者以别的方式想在更一般的背景下理解\trCoherentSheaf{}的人。即便是这些人，也应当等到很久之后再读它。除非你真的想读，否则别读！

设 $A$ 是一个\trRing。回忆一个 $A$-\trModule{} $M$ 何时称为\trFinitelyGenerated{}的、\trFinitelyPresented{}的与\trCoh{}的（\S6.4）。我们喜欢凝聚性的理由是：\trCoh{}\trModule{}构成一个\trAbelianCategory。下面是一些不难入手的练习，把这些概念为何表现良好亲手算清。其中有些重复了先前的讨论，为的是让本节内容自足。

一个\trModule{}的凝聚性这一概念，只有在环于自身之上\trCoh{}的场合才有意思。类似地，概形 $X$ 上的\trCoherentSheaf{}也只有在 $\OO_X$ \trCoh{}（「在自身之上」）时才有意思。在这种情形，凝聚性显然就与有限表现是同一回事。非\trNoetherian{}凝聚性出场的一个例子，是「\trRestrictedPowerSeries{}」（restricted power series）环 $R\langle x_1,\dots,x_n\rangle$，其中 $R$ 是带非离散赋值的 $K$（例如 $\mathbb{Q}_p$ 的\trAlgebraicClosure{}的一个完备化）的\trValuationRing{}。这与\trTate{}的 $K$ 上\trNonArchimedean{}解析几何理论有关（例如可在 [BCDKT] 中读到）。结构层凝聚性的重要性，凸显了复几何中\trOka{}定理（陈述见 \S6.4.8）的重要性。

\textbf{6.7.A.（练习）}\quad 证明：\trCoh{}蕴涵\trFinitelyPresented，而后者又蕴涵\trFinitelyGenerated。（这在 \S6.4 开头讨论过。）

\textbf{6.7.B.（练习）}\quad 证明：$0$ 是\trCoh{}的。

对练习~6.7.C--6.7.I，设
\begin{equation*}\tag{6.7.0.1}\label{eq:coherent-exercise-sequence}
0\longrightarrow M\longrightarrow N\longrightarrow P\longrightarrow 0
\end{equation*}
是 $A$-\trModule{}的一个\trExactSequence。在这一系列练习中，我们将证明：$\{M,N,P\}$ 中若有两个是\trCoh{}的，则第三个亦然——这将会大有用处。

\textbf{6.7.1.（提示）}\quad 下面的提示适用于其中好几道练习：试着写出
\[
\begin{tikzcd}[row sep=2.6em, column sep=2.8em]
0 \ar[r] & A^{\oplus p} \ar[r] \ar[d] & A^{\oplus(p+q)} \ar[r] \ar[d] & A^{\oplus q} \ar[r] \ar[d] & 0 \\
0 \ar[r] & M \ar[r] & N \ar[r] & P \ar[r] & 0{,}
\end{tikzcd}
\]
并可能用上蛇形引理~1.6.5。

\textbf{6.7.C.（练习）}\quad 证明：$N$ \trFinitelyGenerated{}蕴涵 $P$ \trFinitelyGenerated。（你只需要 \eqref{eq:coherent-exercise-sequence} 的\trRightExact{}性。）

\textbf{6.7.D.（练习）}\quad 证明：$M$、$P$ 都\trFinitelyGenerated{}蕴涵 $N$ \trFinitelyGenerated。可能的提示：\S6.7.1。（你只需要 \eqref{eq:coherent-exercise-sequence} 的\trRightExact{}性。）

\textbf{6.7.E.（练习）}\quad 证明：$N$、$P$ 都\trFinitelyGenerated{}并不必然蕴涵 $M$ \trFinitelyGenerated。（提示：若 $I$ 是一个\trIdeal，则有 $0\to I\to A\to A/I\to 0$。）

\textbf{6.7.F.（练习）}\quad 证明：$N$ \trCoh、$M$ \trFinitelyGenerated，蕴涵 $M$ \trCoh。（你只需要 \eqref{eq:coherent-exercise-sequence} 的\trLeftExactness。）

\textbf{6.7.G.（练习）}\quad 证明：$N$、$P$ 都\trCoh{}蕴涵 $M$ \trCoh。凝聚性定义（\S6.4.1）中 (i) 的提示：
\[
\begin{tikzcd}[row sep=2.2em, column sep=2.8em]
 & A^{\oplus q} \ar[dd] \ar[dr] & & \\
 & & A^{\oplus p} \ar[d] \ar[dr] & \\
0 \ar[r] & M \ar[r] \ar[d] & N \ar[r] \ar[d] & P \ar[r] \ar[dr] & 0 \\
 & 0 & 0 & & 0.
\end{tikzcd}
\]
（你只需要 \eqref{eq:coherent-exercise-sequence} 的\trLeftExactness。）

\textbf{6.7.H.（练习）}\quad 证明：$M$ \trFinitelyGenerated{}且 $N$ \trCoh{}蕴涵 $P$ \trCoh。凝聚性定义（\S6.4.1）中 (ii) 的提示：\S6.7.1。

\textbf{6.7.I.（练习）}\quad 证明：$M$、$P$ 都\trCoh{}蕴涵 $N$ \trCoh。（提示：\S6.7.1。）

\textbf{6.7.J.（练习）}\quad 证明：\trCoh{}\trModule{}的有限\trDirectSum{}是\trCoh{}的。

\textbf{6.7.K.（练习）}\quad 设 $M$ \trFinitelyGenerated{}而 $N$ \trCoh。若 $\phi\colon M\to N$ 是任意一个映射，证明：$\operatorname{Im}\phi$ 是\trCoh{}的。

\textbf{6.7.L.（练习）}\quad 证明：\trCoh{}\trModule{}之间映射的\trKernel{}与\trCokernel{}都是\trCoh{}的。

至此，我们已验证：\trCoh{}的 $A$-\trModule{}构成 $A$-\trModule{}的\trCategory{}的一个\trAbelian\trSubcategory。（必须检验的事项：$0$ 应属于此集合；它应对有限和封闭；并且它应对取\trKernel{}与\trCokernel{}封闭。）

\textbf{6.7.M.（练习）}\quad 设 $M$ 与 $N$ 是\trCoh{}\trModule{} $P$ 的两个\trCoh{}\trSubmodule。证明：$M+N$ 与 $M\cap N$ 都是\trCoh{}的。（提示：考虑那个「对的」映射 $M\oplus N\to P$。）

\textbf{6.7.N.（练习）}\quad 证明：若 $A$ 是\trCoh{}的（作为 $A$-\trModule），则\trFinitelyPresented{}\trModule{}都是\trCoh{}的。（当然，若\trFinitelyPresented{}\trModule{}都\trCoh，则 $A$ 必\trCoh，因为 $A$ 就是\trFinitelyPresented{}的！）

\textbf{6.7.O.（练习）}\quad 若 $M$ 是\trFinitelyPresented{}的且 $N$ 是\trCoh{}的，证明：$\Hom(M,N)$ 是\trCoh{}的。（提示：$\Hom$ 对第一个变量是\trLeftExact{}的。）

\textbf{6.7.P.（练习）}\quad 若 $M$ 是\trFinitelyPresented{}的，而 $N$ 是\trCoh{}的，证明：$M\otimes_A N$ 是\trCoh{}的。
